\section{Closed success sets for subspace closeness}

Fix a target subspace $T$ of a finite-dimensional complex coordinate space.
The condition that a candidate subspace is close to $T$ is a closed
condition on its orthogonal projection. This uses the original closeness
condition of \cite[Definition \texttt{d:close}, lines 469--474]
{Arad2016RigorousRG}.

\begin{theorem}[A closed compressed-matrix inequality]
    \label{thm:closed_compressed_positivity}
    \lean{Matrix.isClosed_setOf_compressed_posSemidef}
    \leanok
    Fix a square complex matrix $Q$ and a real number $\delta$. Then
    \begin{align}
        \bigl\{P:QPQ-(1-\delta)Q\ge0\bigr\}
    \end{align}
    is closed in the space of square matrices. No symmetry or projection
    assumption on the fixed matrix $Q$ is needed for this assertion.
\end{theorem}
\begin{proof}\leanok
    The cone of positive-semidefinite matrices is closed. Take its
    inverse image under the continuous affine map
    $P\mapsto QPQ-(1-\delta)Q$.
\end{proof}

\begin{theorem}[Exact matrix form of subspace closeness]
    \label{thm:closeness_compressed_matrix}
    \lean{Submodule.isCloseTo_iff_compressed_posSemidef}
    \leanok
    \uses{def:subspace_close}
    Let $P$ and $Q$ be the actual orthogonal projection matrices onto
    subspaces $U$ and $T$, respectively. Then $U$ is $\delta$-close to $T$
    if and only if $QPQ-(1-\delta)Q$ is positive semidefinite.
\end{theorem}
\begin{proof}\leanok
    The standard matrix representation preserves multiplication,
    subtraction, scalar multiplication and positivity of operators.
    Thus the inequality is exactly the defining condition
    $(1-\delta)P_T\le P_TP_UP_T$.
\end{proof}

\begin{theorem}[The actual success set and its complement]
    \label{thm:closed_closeness_success}
    \lean{Submodule.isClosed_setOf_isCloseTo_projectionRange}
    \leanok
    Let $\mathcal P$ be the space of all orthogonal projection matrices,
    with its inherited matrix topology. For every real $\delta$ and
    every fixed target subspace $T$, the set
    \begin{align}
        \bigl\{P\in\mathcal P:
          \operatorname{range}P\text{ is }\delta\text{-close to }T\bigr\}
    \end{align}
    is closed. Hence this success set and its true failure complement
    are Borel measurable. Zero-rank candidates and the zero target are
    included.
\end{theorem}
\begin{proof}\leanok
    \uses{thm:closed_compressed_positivity,thm:closeness_compressed_matrix,
        thm:projection_matrix_range}
    An orthogonal projection is exactly the orthogonal projection onto
    its actual range. The matrix criterion therefore applies with the
    candidate matrix $P$ itself. Restricting the closed matrix set to
    $\mathcal P$ gives the result.
\end{proof}

No measurable choice of a comparison subspace is used. To apply the result
to an actual random output, one must separately prove that its orthogonal
projection is a measurable function of the sampling variables.
